\begin{abstract}
Machine-learning models for crop-yield prediction are usually evaluated on random splits with
complete weather inputs, although in practice weather records are incomplete and models are
applied to districts that were not seen during training. These two conditions were studied here,
separately and together, for 32 districts of Gujarat, India. Official district statistics for
groundnut, bajra and wheat from 1997--98 to 2022--23 were linked with daily NASA POWER weather, and
18 seasonal features were derived. Ridge regression, a random forest and XGBoost were evaluated
under a random split, a temporal split and a district holdout, and missing weather was simulated by
random masking, removal of variable groups and outages in the daily record at 10 to 40\%, with
median, nearest-neighbour, iterative and native handling. Holding out districts lowered $R^2$ from
0.37 to 0.29 for the random forest, and the random split overstated the accuracy available in
unseen districts with 30\% missing weather by 0.12 to 0.17 in $R^2$. Missing weather raised the
normalised error by up to 1.6 percentage points, most when whole variable groups were missing, and
the choice of model mattered more than the imputer: a random-forest strategy was best in 35 of 36
settings. Contrary to expectation, the two problems did not compound; their combined loss was
significantly smaller than the sum of their separate losses ($p < 0.001$). Predicting future
seasons was the hardest task, since tree ensembles could not extrapolate the yield trend. District
holdouts should therefore replace random splits when yield models are validated for new regions.
\end{abstract}

\begin{IEEEkeywords}
crop yield prediction, missing data, imputation, distribution shift, spatial cross-validation,
random forest, NASA POWER
\end{IEEEkeywords}
